\documentclass[11pt,a4paper]{article}

% --- PAQUETES DE FORMATO Y TIPOGRAFÍA ---
\usepackage[utf8]{inputenc}
\usepackage[T1]{fontenc}
\usepackage[spanish,es-tabla,es-noquoting]{babel}
\usepackage[margin=2.2cm,headheight=15pt]{geometry}
\usepackage[skip=6pt plus 2pt, indent=0pt]{parskip}
\usepackage{amsmath,amssymb,amsfonts}
\usepackage{graphicx}
\usepackage{booktabs}
\usepackage{tabularx}
\usepackage{xcolor}
\usepackage{fancyhdr}
\usepackage{listings}
\usepackage{enumitem}
\usepackage{titlesec}
\usepackage{tcolorbox}
\tcbuselibrary{skins,breakable}
\usepackage{tikz}
\usetikzlibrary{arrows.meta,calc,positioning,shapes.geometric,shapes.misc,fit}
\usepackage{hyperref}

% --- PALETA DE COLORES TÉCNICA Y CORPORATIVA ---
\definecolor{primary}{RGB}{15, 23, 42}       % Slate 900
\definecolor{accent}{RGB}{14, 165, 233}      % Sky 500
\definecolor{secondary}{RGB}{99, 102, 241}   % Indigo 500
\definecolor{success}{RGB}{16, 185, 129}     % Emerald 500
\definecolor{warning}{RGB}{245, 158, 11}     % Amber 500
\definecolor{danger}{RGB}{239, 68, 68}       % Rose 500
\definecolor{lightbg}{RGB}{248, 250, 252}    % Slate 50
\definecolor{codebg}{RGB}{241, 245, 249}     % Slate 100
\definecolor{bordercol}{RGB}{226, 232, 240}  % Slate 200

% --- CONFIGURACIÓN DE HYPERREF ---
\hypersetup{
    colorlinks=true,
    linkcolor=secondary,
    citecolor=primary,
    urlcolor=accent,
    pdftitle={Manual de Usuario: Agent IA},
    pdfauthor={Miguel Angel Polo Castro}
}

% --- CONFIGURACIÓN DE ENCABEZADOS Y PIE DE PÁGINA ---
\pagestyle{fancy}
\fancyhf{}
\fancyhead[L]{\small\color{primary}\textbf{Agent IA} --- Manual de Usuario y Guía de Operación}
\fancyhead[R]{\small\color{gray}\today}
\fancyfoot[C]{\small\color{primary}\thepage}
\renewcommand{\headrulewidth}{0.4pt}
\renewcommand{\footrulewidth}{0.4pt}

% --- FORMATO DE TÍTULOS ---
\titleformat{\section}{\Large\bfseries\color{primary}}{\thesection.}{0.5em}{}[\titlerule]
\titleformat{\subsection}{\large\bfseries\color{secondary}}{\thesubsection.}{0.5em}{}
\titleformat{\subsubsection}{\normalsize\bfseries\color{primary}}{\thesubsubsection.}{0.5em}{}

% --- CAJAS DE INFORMACIÓN Y DESTACADOS (tcolorbox) ---
\newtcolorbox{infobox}[1]{
    colback=lightbg,
    colframe=accent,
    coltitle=white,
    fonttitle=\bfseries\small,
    title=#1,
    arc=2mm,
    boxrule=0.8pt,
    left=3mm, right=3mm, top=2mm, bottom=2mm
}

\newtcolorbox{alertbox}[1]{
    colback=lightbg,
    colframe=danger,
    coltitle=white,
    fonttitle=\bfseries\small,
    title=#1,
    arc=2mm,
    boxrule=0.8pt,
    left=3mm, right=3mm, top=2mm, bottom=2mm
}

\newtcolorbox{tipbox}[1]{
    colback=lightbg,
    colframe=success,
    coltitle=white,
    fonttitle=\bfseries\small,
    title=#1,
    arc=2mm,
    boxrule=0.8pt,
    left=3mm, right=3mm, top=2mm, bottom=2mm
}

% --- CONFIGURACIÓN DE LISTINGS (CÓDIGO) ---
\lstset{
    backgroundcolor=\color{codebg},
    basicstyle=\ttfamily\footnotesize\color{primary},
    keywordstyle=\bfseries\color{secondary},
    stringstyle=\color{success},
    commentstyle=\itshape\color{gray},
    numberstyle=\tiny\color{gray},
    numbers=none,
    frame=single,
    rulecolor=\color{bordercol},
    breaklines=true,
    breakatwhitespace=true,
    showstringspaces=false,
    tabsize=2,
    inputencoding=utf8,
    extendedchars=true,
    literate=
      {á}{{\'a}}1
      {é}{{\'e}}1
      {í}{{\'i}}1
      {ó}{{\'o}}1
      {ú}{{\'u}}1
      {ñ}{{\~n}}1
      {Á}{{\'A}}1
      {É}{{\'E}}1
      {Í}{{\'I}}1
      {Ó}{{\'O}}1
      {Ú}{{\'U}}1
      {Ñ}{{\~N}}1
      {—}{{---}}1
      {¿}{{?`}}1
      {¡}{{!`}}1
}

% --- ESTILOS DE BLOQUES TIKZ ---
\tikzset{
    block/.style={draw=primary, fill=lightbg, thick, rectangle, rounded corners=3pt, align=center, font=\small},
    agentNode/.style={draw=secondary, fill=lightbg, thick, rectangle, rounded corners=3pt, minimum height=2.2em, minimum width=6.5em, align=center, font=\footnotesize\bfseries},
    coreNode/.style={draw=primary, fill=primary!10, very thick, rectangle, rounded corners=4pt, minimum height=2.8em, minimum width=8.5em, align=center, font=\small\bfseries},
    storageNode/.style={draw=accent, fill=accent!10, thick, cylinder, shape border rotate=90, aspect=0.25, minimum height=2.5em, minimum width=6.0em, align=center, font=\footnotesize\bfseries}
}

\begin{document}

% ====================================================================
% PORTADA FORMAL
% ====================================================================
\begin{titlepage}
    \centering
    \vspace*{1.0cm}
    {\scshape\LARGE\bfseries\color{primary} SISTEMA MULTI-AGENTE INTELIGENTE \par}
    \vspace{0.4cm}
    {\scshape\large Copiloto de Conocimiento, Estudio y Segundo Cerebro \par}
    \vspace{1.5cm}
    
    \rule{\linewidth}{1.8pt}\\[0.4cm]
    {\Huge\bfseries\color{primary} Manual Oficial de Usuario\\[0.2cm] y Guía Técnica de Operación \par}
    \vspace{0.2cm}
    {\large\color{accent}\textbf{Agent IA --- Versión 1.0 (Arquitectura Híbrida)} \par}
    \rule{\linewidth}{1.8pt}\\[1.5cm]
    
    \begin{minipage}{0.48\textwidth}
        \begin{flushleft} \large
            \textbf{Autor y Desarrollador:}\\
            Miguel Ángel Polo Castro\\[0.3cm]
            \textbf{Entorno de Ejecución:}\\
            FastAPI + Streamlit + Ollama + Gemini
        \end{flushleft}
    \end{minipage}
    ~
    \begin{minipage}{0.48\textwidth}
        \begin{flushright} \large
            \textbf{Integraciones Principales:}\\
            Notion (29 DBs Relacionales)\\
            Obsidian (Vault Markdown)\\
            ChromaDB (Memoria Vectorial)
        \end{flushright}
    \end{minipage}
    
    \vfill
    
    \begin{infobox}{Resumen del Documento}
        Esta guía detalla la configuración, despliegue, operación y buenas prácticas del sistema \textbf{Agent IA}. Se documenta la interacción en lenguaje natural con el orquestador \textbf{Hermes}, el funcionamiento de la flota de agentes especializados (\textbf{Curador}, \textbf{Plan}, \textbf{Estudio}, \textbf{Sync}), el sistema de protección y ofuscación de datos locales (\texttt{DataMasker}) y el modo híbrido conmutable entre modelos locales y en la nube.
    \end{infobox}
    
    \vspace{0.8cm}
    {\large \today \par}
\end{titlepage}

\tableofcontents
\newpage

% ====================================================================
\section{Introducción y Visión General}
% ====================================================================

\subsection{¿Qué es Agent IA?}
\textbf{Agent IA} es un sistema multi-agente personal diseñado para operar como un copiloto proactivo de gestión de conocimiento, estudio y planificación estratégica. El sistema conecta y mantiene sincronizados los dos pilares del \textbf{Segundo Cerebro}:
\begin{enumerate}[leftmargin=1.5cm]
    \item \textbf{Notion:} Base relacional de alta estructuración con 29 bases de datos interconectadas (Objetivos, Proyectos, Tareas, Áreas, Recursos, etc.).
    \item \textbf{Obsidian:} Vault local de notas en formato Markdown plano, garantizando portabilidad, velocidad y control total de archivos.
\end{enumerate}

Toda la interacción con el usuario se realiza mediante \textbf{lenguaje natural} a través de \textbf{Hermes}, un agente orquestador central que clasifica intenciones, propone acciones y delega la ejecución en agentes especializados.

\subsection{Filosofía de Diseño: Human-in-the-Loop y Revisión Proactiva}
Agent IA opera bajo el principio estricto de \textbf{revisión humana}:
\begin{itemize}
    \item \textbf{Propone, no destruye:} Las notas, metas o tareas nunca se escriben de forma destructiva sin confirmación previa si implican decisiones de categorización o estructuración compleja.
    \item \textbf{Aprobación conversacional:} El usuario puede confirmar propuestas mediante lenguaje natural (\textit{"sí", "dale", "hazlo", "perfecto"}) o a través de tarjetas interactivas con botones en la interfaz web.
\end{itemize}

% ====================================================================
\section{Arquitectura del Sistema y Flota de Agentes}
% ====================================================================

Agent IA está construido bajo una \textbf{Arquitectura Hexagonal (Ports \& Adapters)}, lo que desacopla la lógica de negocio y agentes de los proveedores de inteligencia artificial (Ollama, Gemini) y de almacenamiento (Obsidian, Notion, ChromaDB).

\begin{figure}[!htbp]
\centering
\begin{tikzpicture}[auto, node distance=2.2cm, >=Latex]
    % Frontend
    \node [coreNode, fill=secondary!15] (ui) {Frontend Web AI-Native\\(\texttt{:8000}) / HUD (\texttt{:8501})};
    
    % Gateway / Hermes
    \node [coreNode, below of=ui, node distance=2.0cm] (hermes) {\textbf{Hermes (Orquestador)}\\Clasificador Semántico \& Máquina de Estados};
    
    % Flota de Agentes
    \node [agentNode, below of=hermes, xshift=-4.5cm, node distance=2.2cm] (curador) {AgenteCurador\\(Notas \& Tags)};
    \node [agentNode, below of=hermes, xshift=-1.5cm, node distance=2.2cm] (plan) {AgentePlan\\(Metas \& Tareas)};
    \node [agentNode, below of=hermes, xshift=1.5cm, node distance=2.2cm] (estudio) {AgenteEstudio\\(Repaso SM-2)};
    \node [agentNode, below of=hermes, xshift=4.5cm, node distance=2.2cm] (sync) {AgenteSync\\(Sincronización)};
    
    % Almacenamiento e Integraciones
    \node [storageNode, below of=curador, node distance=2.2cm] (obsidian) {Obsidian Vault\\(Local MD)};
    \node [storageNode, below of=plan, node distance=2.2cm] (notion) {Notion DBs\\(29 Tablas)};
    \node [storageNode, below of=estudio, node distance=2.2cm] (chroma) {ChromaDB\\(Embeddings)};
    \node [storageNode, below of=sync, node distance=2.2cm] (llm) {LLM Engines\\(Gemini / Ollama)};

    % Conexiones
    \draw [<->, thick] (ui) -- (hermes);
    \draw [->, thick] (hermes) -- (curador);
    \draw [->, thick] (hermes) -- (plan);
    \draw [->, thick] (hermes) -- (estudio);
    \draw [->, thick] (hermes) -- (sync);

    \draw [<->, dashed] (curador) -- (obsidian);
    \draw [<->, dashed] (curador) -- (notion);
    \draw [<->, dashed] (plan) -- (notion);
    \draw [<->, dashed] (curador) -- (chroma);
    \draw [<->, dashed] (hermes) -- (llm);
\end{tikzpicture}
\caption{Diagrama de arquitectura hexagonal y flujo de delegación multi-agente.}
\label{fig:arch_diagram}
\end{figure}

\subsection{Descripción de los Agentes Especializados}

\begin{table}[!htbp]
\centering
\caption{Roles y responsabilidades de la flota de agentes.}
\label{tab:agentes_summary}
\begin{tabularx}{\textwidth}{@{}llX@{}}
\toprule
\textbf{Agente} & \textbf{Dominio} & \textbf{Responsabilidad Principal} \\ \midrule
\textbf{Hermes} & Orquestación & Analiza el prompt, gestiona memoria de sesión, rutea peticiones y administra el ciclo de confirmación activa. \\
\textbf{AgenteCurador} & Nota, Área & Extrae taxonomía (áreas, temas, tags), detecta duplicados vía búsqueda semántica y persiste en Notion + Obsidian. \\
\textbf{AgentePlan} & Tarea, Proyecto, Meta & Desglosa metas en planes estructurados (Objetivos $\rightarrow$ Proyectos $\rightarrow$ Tareas) y ejecuta inserciones relacionales. \\
\textbf{AgenteEstudio} & ItemEstudio & Administra tarjetas de estudio y calcula intervalos de repetición espaciada con el algoritmo SuperMemo SM-2. \\
\textbf{AgenteSync} & Integraciones & Vela por la consistencia de datos entre plataformas externas (Notion, Obsidian, Todoist, Calendar). \\ \bottomrule
\end{tabularx}
\end{table}

% ====================================================================
\section{Requisitos y Preparación del Entorno}
% ====================================================================

\subsection{Prerrequisitos de Software}
\begin{itemize}
    \item \textbf{Sistema Operativo:} Windows 10/11 o Linux (Arch, Ubuntu, Debian).
    \item \textbf{Python:} Versión 3.14 o superior.
    \item \textbf{Gestor de Paquetes:} \texttt{uv} (Astral) para sincronización determinista de dependencias.
    \item \textbf{Ollama (Local):} Para generación de embeddings locales (\texttt{nomic-embed-text}) y fallback offline.
    \item \textbf{Google AI Studio API Key (Opcional pero Recomendado):} Para inferencia instantánea sub-segundo con Gemini Flash.
\end{itemize}

\subsection{Instalación de Dependencias con \texttt{uv}}
Desde la raíz del repositorio, ejecute:
\begin{lstlisting}[language=bash]
# Sincronizar el entorno virtual completo
uv sync
\end{lstlisting}

\subsection{Modelos Locales de Ollama}
Asegúrese de descargar los pesos mínimos necesarios:
\begin{lstlisting}[language=bash]
# Modelo ultraligero de embeddings (Memoria Vectorial ChromaDB)
ollama pull nomic-embed-text

# Modelo local ligero para fallback / offline
ollama pull llama3.2:3b
\end{lstlisting}

% ====================================================================
\section{Configuración de Variables de Entorno (\texttt{.env})}
% ====================================================================

El archivo \texttt{.env} controla la conectividad, las claves criptográficas y el modo de operación del motor de inferencia.

\begin{lstlisting}[caption={Estructura completa de configuración en .env}]
# ===============================================================
# Agent IA — Configuración Centralizada
# ===============================================================

# --- Integracion Notion ---
AGENT_NOTION_TOKEN=ntn_xxxxxxxxxxxxxxxxxxxxxxxxxxxxxxxxxxxxxxxxxxx
AGENT_NOTION_ROOT_PAGE_ID=6941a7f7-a5d2-463f-a718-6cb9030f7c8c

# --- Motor Local Ollama ---
AGENT_OLLAMA_BASE_URL=http://localhost:11434
AGENT_OLLAMA_MODEL=llama3.2:3b
AGENT_OLLAMA_MODEL_RAPIDO=llama3.2:3b
AGENT_OLLAMA_EMBED_MODEL=nomic-embed-text
AGENT_OLLAMA_KEEP_ALIVE=30m

# --- Almacenamiento Local (Obsidian y ChromaDB) ---
AGENT_OBSIDIAN_VAULT_PATH=C:\Users\migue\Agent_IA_Vault
AGENT_CHROMA_PERSIST_DIR=./data/chroma

# --- Backend LLM e Inferencia en la Nube ---
AGENT_LLM_BACKEND=hybrid                 # Opciones: hybrid | gemini | ollama
AGENT_GEMINI_API_KEY=AIzaSyxxxxxxxxxxxx  # Clave de Google AI Studio
AGENT_GEMINI_MODEL=gemini-2.0-flash
AGENT_GEMINI_MODEL_PLAN=gemini-2.0-flash
AGENT_ENABLE_DATA_MASKING=true           # Anonimizacion local de privacidad

# --- Puertos de Red ---
AGENT_API_HOST=0.0.0.0
AGENT_API_PORT=8000
AGENT_HUD_PORT=8501
\end{lstlisting}

\begin{infobox}{Modos de Backend (\texttt{AGENT\_LLM\_BACKEND})}
\begin{itemize}
    \item \textbf{hybrid (Recomendado):} Hermes, Curador y Plan utilizan Gemini Flash ($<800\text{ ms}$), mientras que ChromaDB y los embeddings permanecen 100\% locales con \texttt{nomic-embed-text}.
    \item \textbf{ollama:} Inferencia 100\% local en tu máquina (sin conexión a internet).
    \item \textbf{gemini:} Todo el pipeline opera en la nube de Google.
\end{itemize}
\end{infobox}

% ====================================================================
\section{Puesta en Marcha y Servicios}
% ====================================================================

El sistema se compone de dos interfaces independientes servidas por FastAPI y Streamlit:

\begin{figure}[!htbp]
\centering
\begin{minipage}{0.48\textwidth}
    \begin{lstlisting}[language=bash, caption={Terminal 1: API Gateway}]
# Iniciar Servidor FastAPI y Web
uv run agent-api
    \end{lstlisting}
\end{minipage}
~
\begin{minipage}{0.48\textwidth}
    \begin{lstlisting}[language=bash, caption={Terminal 2: HUD Streamlit (Opcional)}]
# Iniciar Panel Analitico Secundario
uv run agent-hud
    \end{lstlisting}
\end{minipage}
\end{figure}

\subsection{Accesos Disponibles en el Navegador}
\begin{itemize}
    \item \textbf{Interfaz Web AI-Native Principal:} \href{http://localhost:8000/}{http://localhost:8000/}
    \item \textbf{Documentación Interactiva Swagger / OpenAPI:} \href{http://localhost:8000/docs}{http://localhost:8000/docs}
    \item \textbf{HUD de Métricas en Streamlit:} \href{http://localhost:8501/}{http://localhost:8501/}
\end{itemize}

% ====================================================================
\section{Guía de Interacción y Casos de Uso}
% ====================================================================

\subsection{1. Curaduría y Captura de Conocimiento (\texttt{AgenteCurador})}
Para guardar una nota, recurso o apunte, basta con indicarlo a Hermes:

\begin{infobox}{Ejemplo de Prompt para Curaduría}
\textbf{Usuario:} \textit{"Anota esto: Protocolo Raft resuelve consenso distribuido mediante elección de líder y replicación de log garantizando consistencia fuerte."}
\end{infobox}

\textbf{Flujo de Ejecución:}
\begin{enumerate}
    \item Hermes detecta la intención y delega a \texttt{AgenteCurador}.
    \item \texttt{AgenteCurador} genera el embedding del texto y busca notas similares en ChromaDB para prevenir duplicados.
    \item El LLM categoriza la nota:
    \begin{itemize}
        \item \textbf{Área sugerida:} \texttt{Sistemas Distribuidos} / \texttt{Redes}
        \item \textbf{Tags:} \texttt{consenso}, \texttt{raft}, \texttt{distribuidos}
    \end{itemize}
    \item El sistema presenta una tarjeta interactiva de confirmación.
    \item Al responder \textbf{\textit{"sí"}} o presionar el botón \textbf{\texttt{[ Confirmar ]}}:
    \begin{itemize}
        \item Se crea la página relacional en Notion con tags y área vinculada.
        \item Se genera el archivo \texttt{Redes/Protocolo Raft.md} en el Vault de Obsidian.
        \item Se indexa el ID de Notion en ChromaDB para búsquedas futuras.
    \end{itemize}
\end{enumerate}

\subsection{2. Planificación Estratégica de Proyectos y Metas (\texttt{AgentePlan})}

\begin{infobox}{Ejemplo de Prompt para Planificación}
\textbf{Usuario:} \textit{"Quiero un plan estratégico para preparar mi examen final de Redes de Computadores en 3 semanas."}
\end{infobox}

\textbf{Resultado Estructurado que Genera el Sistema:}
\begin{itemize}
    \item \textbf{Objetivo Principal:} \textit{Aprobar Examen Final de Redes de Computadores} (Área: \textit{Universidad}).
    \item \textbf{Proyecto 1: Dominio de Capa Física y Enlace}
    \begin{itemize}
        \item $[~]$ Resolver laboratorio de modulación y sincronismo.
        \item $[~]$ Elaborar tabla comparativa de tramas UART vs SPI.
    \end{itemize}
    \item \textbf{Proyecto 2: Enrutamiento y Capa de Red}
    \begin{itemize}
        \item $[~]$ Configurar topología OSPF y BGP en simulador.
        \item $[~]$ Practicar ejercicios de subnetting VLSM.
    \end{itemize}
\end{itemize}

Al confirmar, el sistema inserta automáticamente las entidades en cascada en las bases de datos de Notion (\texttt{Objetivos} $\rightarrow$ \texttt{Proyectos} $\rightarrow$ \texttt{Tareas}) estableciendo las relaciones bidireccionales.

\subsection{3. Creación Rápida de Tareas (Fast-Path)}
Si solo necesitas agregar una tarea atómica sin desglose estratégico:
\begin{infobox}{Comando Rápido}
\textbf{Usuario:} \textit{"agrega una tarea: Enviar informe de laboratorio a David Herrera antes de las 6pm"}
\end{infobox}
El sistema omite el procesamiento de LLM, ejecuta la inserción directa en la base de datos de \texttt{Tareas} de Notion y responde en menos de 300 ms.

\subsection{4. Study Board y Repetición Espaciada SM-2 (\texttt{AgenteEstudio})}
El módulo de estudio activo permite consolidar conceptos a largo plazo:

\begin{infobox}{Ejemplo de Generación de Flashcards}
\textbf{Usuario:} \textit{"Genera flashcards sobre el algoritmo de consenso Raft"}
\end{infobox}

\textbf{Flujo Operativo de Estudio:}
\begin{itemize}
    \item \textbf{Extracción Atómica (RF3.1):} Sintetiza pares de Pregunta y Respuesta concisos y los persiste en \texttt{data/study\_board.json}.
    \item \textbf{Notas Cornell (RF3.3):} Con el comando \textit{"Crea una nota Cornell sobre X"}, estructura apuntes en Obsidian con ideas clave, notas y resumen.
    \item \textbf{Sesión Interactiva Quiz SM-2 (RF3.4):} Al escribir \textit{"Iniciar repaso"}, el agente presenta tarjetas una por una. Al calificar de 0 a 5, el algoritmo recalibra el factor de facilidad e intervalo en días.
\end{itemize}

\subsection{5. Sincronización Proactiva y Webhooks (\texttt{AgenteSync})}
Garantiza la consistencia continua entre las distintas herramientas del ecosistema:

\begin{infobox}{Ejemplo de Sincronización y Alertas}
\textbf{Usuario:} \textit{"Sincroniza mis notas de Notion con Obsidian y dime qué tareas vencen hoy"}
\end{infobox}

\textbf{Capacidades del AgenteSync (Fase 3):}
\begin{itemize}
    \item \textbf{Sincronización Bidireccional (RF1.1):} Rastrea hashes SHA-256 en \texttt{data/sync\_state.json} e indexa cambios en ChromaDB.
    \item \textbf{Alertas Proactivas de Tareas (RF6.3):} Identifica tareas vencidas o con entrega en las próximas 48 horas.
    \item \textbf{Gateway para Webhooks n8n:} Endpoints HTTP \texttt{POST /sync/ejecutar} y \texttt{POST /webhook/sync} para recibir eventos automáticos desde Telegram, Todoist o n8n.
\end{itemize}

% ====================================================================
\section{Visualizador 3D Isométrico ("The Sims for AI Agents")}
% ====================================================================

El sistema incorpora un entorno espacial 3D interactivo accesible en \href{http://localhost:8000/sims}{http://localhost:8000/sims} o mediante el botón \textbf{\texttt{[ Oficina 3D | SIMS ]}} en la barra superior.

\subsection{Estaciones de Trabajo y Gemelos Digitales}
\begin{itemize}
    \item \textbf{Desk 01 (Active Compute):} Puesto de mando ultrawide donde opera el avatar de \textbf{Hermes (Azul)}.
    \item \textbf{The Vault (Secure Storage):} Bóveda acorazada donde el avatar de \textbf{AgenteCurador (Naranja)} resguarda notas y ChromaDB.
    \item \textbf{Kanban Wall (Work Items):} Pizarra de cristal donde el avatar de \textbf{AgentePlan (Morado)} gestiona metas en cascada.
    \item \textbf{Study Board (SM-2 Lab):} Biblioteca holográfica donde el avatar de \textbf{AgenteEstudio (Verde)} procesa flashcards.
    \item \textbf{Security Gate (Access Control):} Pasarela óptica donde el avatar de \textbf{AgenteSync (Amarillo)} supervisa el DataMasker.
\end{itemize}

\subsection{Controles del Simulador 3D}
\begin{itemize}
    \item \textbf{Rotación y Zoom:} Clic izquierdo sostenido para orbitar en 360° y rueda del ratón para zoom.
    \item \textbf{Inspección de Telemetría:} Al hacer clic sobre cualquier agente o estación se despliega la tarjeta inferior con \textit{Runtime}, \textit{Tokens consumidos}, \textit{Estado de ejecución} y barra de progreso.
    \item \textbf{Modo Paseo e Iluminación:} Botones para activar el desplazamiento autónomo de los robots y alternar entre modo Día y Noche (\textit{Cyberpunk Dark-Glow}).
\end{itemize}

% ====================================================================
\section{Sistema de Privacidad y Anonimización (\texttt{DataMasker})}
% ====================================================================

Para garantizar la confidencialidad de la información cuando se utiliza el backend de Gemini Flash en la nube, el sistema incorpora el middleware \texttt{DataMasker} (\texttt{sanitizer.py}).

\subsection{Mecanismo de Enmascaramiento y Restitución Local}
\begin{figure}[!htbp]
\centering
\begin{tikzpicture}[node distance=1.8cm, auto, >=Latex]
    \node [block, text width=9.5cm, fill=lightbg] (in) {\textbf{Texto Original (Local):}\\ \textit{"Contactar a soporte@empresa.com con api\_key=AIzaSy9843hf por \$5,000 USD"}};
    \node [block, text width=9.5cm, fill=warning!10, below of=in] (mask) {\textbf{DataMasker (Filtro Regex Local):}\\ Enmascara PII y genera identificadores sintéticos};
    \node [block, text width=9.5cm, fill=codebg, below of=mask] (out) {\textbf{Payload Transmitido a Gemini por HTTPS:}\\ \textit{"Contactar a <EMAIL\_1> con <SECRET\_1> por <AMOUNT\_1>"}};
    \node [block, text width=9.5cm, fill=success!10, below of=out] (resp) {\textbf{Desofuscador Local (Post-Respuesta):}\\ Restituye los valores originales antes de guardar en Notion/Obsidian};

    \draw [->, thick] (in) -- (mask);
    \draw [->, thick] (mask) -- (out);
    \draw [->, thick] (out) -- (resp);
\end{tikzpicture}
\caption{Flujo de sanitización local y protección de privacidad.}
\label{fig:datamasker_flow}
\end{figure}

\subsection{Entidades Protegidas Automáticamente}
\begin{itemize}
    \item \textbf{Secretos y Credenciales:} \texttt{api\_key}, \texttt{token}, \texttt{password}, \texttt{contraseña}.
    \item \textbf{Direcciones de Correo:} Cualquier formato estándar \texttt{user@domain.ext}.
    \item \textbf{Números Telefónicos:} Secuencias locales e internacionales de 7 a 15 dígitos.
    \item \textbf{Montos Financieros:} Cifras con símbolos de moneda (\$, €, USD, COP, EUR, etc.).
\end{itemize}

% ====================================================================
\section{Diagnóstico y Resolución de Problemas}
% ====================================================================

\begin{table}[!htbp]
\centering
\caption{Guía de diagnóstico de errores comunes.}
\label{tab:troubleshooting}
\begin{tabularx}{\textwidth}{@{}lXX@{}}
\toprule
\textbf{Síntoma} & \textbf{Causa Probable} & \textbf{Solución} \\ \midrule
\textbf{Error 429 Too Many Requests} & Límite de peticiones de Google AI Studio alcanzado. & El sistema activa el fallback local automáticamente. Si persiste, espere 60 segundos o configure una API key con facturación. \\
\textbf{Respuestas tardan $>30\text{ s}$} & El backend está en modo \texttt{ollama} ejecutando modelos pesados en CPU. & Cambie en \texttt{.env} a \texttt{AGENT\_LLM\_BACKEND=hybrid} o baje el modelo local a \texttt{llama3.2:3b}. \\
\textbf{Error de conexión en Notion} & Token de Notion expirado o página raíz no compartida con la integración. & Verifique en Notion que la integración \textbf{Agent IA} tenga permisos de lectura/escritura en la página raíz. \\
\textbf{ChromaDB no inicia} & Permisos de escritura bloqueados en la carpeta \texttt{./data/chroma}. & Asegúrese de que el usuario tenga permisos de escritura en la carpeta del proyecto. \\ \bottomrule
\end{tabularx}
\end{table}

% ====================================================================
\section{Comandos de Referencia Rápida}
% ====================================================================

\begin{table}[!htbp]
\centering
\caption{Comandos esenciales de la CLI del proyecto.}
\label{tab:cli_commands}
\begin{tabularx}{\textwidth}{@{}lX@{}}
\toprule
\textbf{Comando} & \textbf{Acción} \\ \midrule
\texttt{uv run agent-api} & Inicia el API Gateway de FastAPI y la interfaz web en \texttt{http://localhost:8000}. \\
\texttt{uv run agent-hud} & Inicia el HUD de monitoreo en Streamlit en \texttt{http://localhost:8501}. \\
\texttt{uv run pytest -v tests/} & Ejecuta la suite completa de 53 pruebas unitarias y de integración. \\
\texttt{python scripts/sync\_notion\_map.py} & Re-escanea las 29 bases de datos de Notion y regenera la caché local. \\
\texttt{ollama list} & Muestra los modelos de inteligencia artificial instalados localmente. \\ \bottomrule
\end{tabularx}
\end{table}

\end{document}
